
@@M0@@
\uline{où on pose}, \uline{pour} @@M1@@ :
@@M2@@
\uline{et pour} @@M3@@ \uline{et} @@M4@@ :
@@M5@@
\uline{les notations étant celles de} 3.1 \uline{et} 3.3, \uline{et} @@M6@@ \uline{désignant la} @@M7@@-\uline{ième classe de Chern de} @@M8@@ (V 6.7). \uline{En outre}, @@M9@@ \uline{admet un inverse à gauche} @@M10@@, \uline{défini par} @@M11@@.

\uline{Si de plus} @@M12@@ \uline{est noethérien et} @@M13@@ \uline{régulier}, \uline{on a la suite exacte} :
@@M14@@

\emph{Démonstration.} Notons qu'on pourrait donner du cas @@M15@@ régulier une démonstration purement algébrique, en termes de théorie des @@M16@@-anneaux, utilisant la description 4.7 de la situation. On va au lieu de cela calquer la démonstration de 3.7.

On sait que les homomorphismes « image inverse » @@M17@@, étant des @@M18@@-homomorphismes, sont compatibles aux filtrations de @@M19@@-anneaux, et par suite ils définissent des homomorphismes sur les gradués associés, d'où les homomorphismes :
@@M20@@
et les homomorphismes analogues obtenus après tensorisation par @@M21@@. Par ailleurs, 4.6 montre qu'après tensorisation par @@M22@@, les homomorphismes @@M23@@ passent aux gradués associés, et définissent
donc :
@@M24@@
